\section*{Fonts}

Per al desenvolupament del projecte s'han consultat diferents tipus de fonts relacionades amb l'anàlisi estàtica de software, els compiladors i les representacions intermèdies de codi.

Com a referència general sobre els processos de compilació i les diferents fases d'anàlisi del codi font, s'ha consultat el llibre \textit{Compilers: Principles, Techniques, and Tools}, d'Aho et al. \cite{aho_compilers}. Aquesta obra proporciona una base teòrica sobre conceptes com l'anàlisi lèxica i sintàctica, l'anàlisi semàntica, les representacions intermèdies i l'anàlisi del flux de control.

També s'ha consultat el treball de Yamaguchi et al. sobre els \textit{Code Property Graphs} \cite{yamaguchi_cpg}, especialment rellevant per al projecte, ja que presenta una representació que combina informació sintàctica, de flux de control i de dependències del programa. Aquesta referència permet situar la proposta del projecte dins del context d'altres aproximacions que utilitzen una representació comuna per a l'anàlisi de codi.

Pel que fa a les eines existents, s'ha consultat la documentació de Tree-sitter \cite{tree_sitter}, Joern \cite{joern_cpg}, CodeQL \cite{codeql} i Clang/LLVM \cite{clang_llvm}. Aquestes eines permeten estudiar diferents aproximacions actuals al parsing, la representació i l'anàlisi de codi, així com les possibilitats d'utilitzar informació comuna entre diferents llenguatges.

També s'han consultat propostes relacionades amb l'abstracció de conceptes de codi, com Koncepto \cite{koncepto}, així com recursos orientats a l'anàlisi estàtica en l'àmbit industrial, com la guia de Static Code Analysis de JetBrains \cite{static_code_analysis_jetbrains}.

Finalment, s'han utilitzat alguns recursos de caràcter més divulgatiu per complementar la consulta de conceptes concrets, com les explicacions sobre arbres sintàctics abstractes \cite{balapriya_ast}, grafs de dependències \cite{gfg_dependencygraph} i grafs de flux de control \cite{gfg_controlflowgraph}.

\printbibliography